\documentclass[11pt,a4paper]{article}
\usepackage[margin=25mm]{geometry}
\usepackage{amsmath,amssymb,amsthm,hyperref}
\setlength{\emergencystretch}{2em}
\newtheorem{theorem}{Theorem}
\newtheorem{lemma}[theorem]{Lemma}
\newcommand{\gen}[1]{\langle#1\rangle}
\newcommand{\Sym}{\operatorname{Sym}}
\newcommand{\Alt}{\operatorname{Alt}}
\title{Faithful augmentation at a symmetric wing}
\author{Human Compression project: a component for L019}
\date{}
\begin{document}\maketitle
Retain the actual perfect-split setting of the acquired theorem
\emph{Symmetric wings at a perfect split}. Thus a transitive string
C-marking $(\rho_0,\ldots,\rho_{r-1})$ acts on $A\sqcup B$,
its left marks are supported on $A$, its right marks on $B$, and
\[
 q=\rho_i=(a\ b)\alpha\beta,
 \quad a\in A,\quad b\in B,
\]
with $\alpha$ supported on $A\setminus\{a\}$ and $\beta$ on
$B\setminus\{b\}$. The omission of $q$ has exactly the two
orbits $A,B$. Write
\[
 L=\gen{\rho_0,\ldots,\rho_{i-1}},\quad
 R=\gen{\rho_{i+1},\ldots,\rho_{r-1}},\quad
 H=\gen{L,q},\quad J=\gen{q,R}.
\]
The groups $L,R$ are actual transitive groups on their supports.
The intersection property gives $L\cap J=1$ and $H\cap J=\gen q$.
We use neither a general quotient intersection property nor a
converse to the endpoint central-twist construction.

\begin{lemma}\label{faithful}
Suppose $L=\Sym(A)$ and $\alpha,\beta$ are both nonidentity.
Restriction of $J$ to $\{a\}\cup B$ is faithful. Its image
\[
 K=\gen{q',R},\qquad q'=(a\ b)\beta,
\]
is a doubly transitive string C-group. It admits a surjective
homomorphism $\chi:K\longrightarrow C_2$ which kills $R$ and
sends $q'$ to the nonidentity element. Moreover
\[
 H=\Sym(A\cup\{b\})\times\gen\beta,
 \qquad K\ne\Alt(\{a\}\cup B),\Sym(\{a\}\cup B).
\]
\end{lemma}
\begin{proof}
Nontriviality of the internal factors implies $|A|,|B|\ge3$.
Since $\alpha\in L$, the group $H$ contains
$q\alpha=(a\ b)\beta$. Choose $c\in A\setminus\{a\}$.
Then $(a\ c)\in L$, and
\[
 \bigl((q\alpha)(a\ c)\bigr)^3=\beta:
\]
the factors on $\{a,b,c\}$ and the support of $\beta$ are
disjoint, and the former is a three-cycle. Thus $H$ contains
$\beta$ and $(a\ b)$. Together with $L$, this transposition
generates $\Sym(A\cup\{b\})$. The reverse inclusion follows
from the displayed supports of the generators, proving the
direct-product formula.

On $A\setminus\{a\}$ every element of $J$ acts as either $1$
or $\alpha$. Consequently the kernel of its restriction to
$\{a\}\cup B$ is contained in $\gen\alpha\le L$.
Since $L\cap J=1$, that kernel is trivial. The actual submarking
of $J$ therefore transfers by an isomorphism to the displayed
marking of $K$, preserving every subgroup intersection and every
nonidentity generator. The group $K$ is transitive, and its
point stabilizer at $a$ contains the transitive group $R$ on
the complement $B$. Hence it is doubly transitive and, in
particular, primitive.

Use the inverse of this restriction isomorphism, followed by
restriction to $A\setminus\{a\}$. This defines the stated
homomorphism $\chi$ onto $\gen\alpha\cong C_2$.
An alternating group has no nontrivial homomorphism to $C_2$,
since it is generated by three-cycles. This excludes that image.
If $K$ were the full symmetric group, its nontrivial character
$\chi$ would be the sign character. The condition $\chi(q')\ne1$
would then make $\beta$ even, because $q'=(a\ b)\beta$.
Thus the unique lift to $J$ of the permutation $\beta\in K$
would act trivially on $A\setminus\{a\}$ and would be precisely
the supported permutation $\beta$. The direct-product formula
would place this nonidentity element in $H\cap J=\gen q$.
But $\beta$ fixes $a$ while $q$ moves $a$, a contradiction.
\end{proof}

\begin{theorem}\label{descent}
Suppose $2r>|A|+|B|+2$, $L=\Sym(A)$, and both internal
factors are nontrivial. Then either
\[
                         2i>|A|+2,                    \tag{1}
\]
or $|B|=5$, $R$ has rank two and order ten, $|A|$ is even,
and $i=|A|/2+1$. In the latter case the faithful image $K$
from Lemma~\ref{faithful} is the degree-six copy of $S_5$,
and its ordered marking, with $q'$ first, is conjugate to one
of the following two tuples:
\[
\begin{split}
 &\bigl((1\ 2)(3\ 4)(5\ 6),\ (2\ 3)(4\ 5),\ (3\ 5)(4\ 6)\bigr),\\
 &\bigl((1\ 2)(3\ 6)(4\ 5),\ (2\ 3)(4\ 5),\ (3\ 5)(4\ 6)\bigr).
\end{split}                                             \tag{2}
\]
Their types are respectively $[6,5]$ and $[4,5]$, and their
attachment point is $a=1$, with $b=2$.
\end{theorem}
\begin{proof}
Put $t=r-i-1$, $m=|B|+1$ and $d=t+1=\operatorname{rank}(K)$.
The exact acquired primitive theorem of L047 applies to $K$:
Lemma~\ref{faithful} supplied its faithful primitive C-action
and excluded both natural alternating and symmetric groups.
If $d<m/2$, integrality gives $2t\le|B|-2$. Consequently
\[
 2i=2r-2-2t>|A|+|B|-2t\ge |A|+2,
\]
as required.

The complete primitive exceptions with $d\ge m/2$ are the
retained degree-six $A_5,S_5$ tuples and the degree-ten rank-five
$S_6$ tuple. We test their twelve exact ordered census
representatives in both orientations. For each tuple we test
every possible anchor, transitivity of its tail on the
complement, and the existence of the character which sends
the first mark to $1\in\mathbb F_2$ and the others to zero.
Only four oriented records survive: source indices
$(3,1),(4,1),(6,0),(8,0)$, where the second coordinate specifies
reversal and indices start at zero. The first and third are
identical and give the first tuple in (2). The second is the
second displayed tuple. The fourth is its conjugate under
the point map $(1,2,3,4,5,6)\mapsto(1,2,4,5,6,3)$.

For clarity, the character test is exact: during finite group
closure attach to each visited permutation the sum modulo two
of its first-generator occurrences. A conflict excludes the
character; if no conflict occurs, every multiplication by
every generator respects the assigned values, which defines
the required homomorphism. Endpoint anchors and complementary
orbits are finite point tests. The independent certificate
retains every tested tuple, the character results and the
explicit conjugacies. Its completeness premise is the acquired
twelve-tuple classification, not this filtering calculation.

In either surviving case $|B|=5$ and $t=2$. The assumed rank
inequality becomes $2i>|A|+1$. If (1) fails, the even integer
$2i$ must equal $|A|+2$. This forces $|A|$ to be even and gives
the stated residual value of $i$. The tail in both tuples has
order ten and is transitive on the five points other than $1$.
\end{proof}

\begin{lemma}\label{inherit}
Suppose the actual left wing has a perfect split at an interior
index $j$, with $1\le j\le i-2$. Then it is also a perfect split
of the full marking. If one of its internal factors was trivial
in the wing, that factor remains trivial in the full marking.
\end{lemma}
\begin{proof}
Write $A=C\sqcup D$ for this wing split, whose crossing pair is
$(c,d)\in C\times D$. The marks before $j$ generate a transitive
group $P$ supported on $C$, and those strictly between $j$ and
$i$ act transitively on $D$. The former list is nonempty, so
$|C|>1$. The split mark $q=\rho_i$ commutes with every member
of $P$; as $C$ is its only nontrivial orbit, $q$ preserves $C$.
It also commutes with $\rho_j$. The latter has exactly one
edge from $C$ to its complement, namely $(c,d)$. Therefore $q$
fixes both $c$ and $d$. An element centralizing a transitive
group and fixing one point fixes the entire orbit, so $q$
fixes $C$ pointwise.

The old attachment point $a$ lies in $D$: the mark
$\rho_{i-1}$ moves $a$ and is supported there, since $j<i-1$.
It follows that the right side of the new omission is transitive
on $D\cup B$, using the old transitive action on $D$, the
crossing $(a,b)$ and the transitive action of $R$ on $B$.
All right-side marks fix $C$. Thus $C,D\cup B$ are precisely
the required two omission orbits, and $\rho_j$ still has its
single crossing edge and the required supported factors.
Those factors have not changed, proving the final assertion.
\end{proof}

\paragraph{Role in the recovery and attribution.}
The separately reviewed actual-wing theorem ensures that at
least one wing is symmetric in the stated high-rank range;
reverse the marking to call it $L$. The present component then
supplies a legitimate restriction isomorphism and the complete
two-type boundary to descending inside that symmetric wing.
The rank theorem and finite primitive classification are the
historical results retained through the complete L047 recovery.
The character and intersection arguments above are new proof
interfaces in this recovery. Lemma~\ref{inherit} is the actual
support argument underlying the corresponding inheritance step
in the source; no priority for that elementary observation is
claimed.

The residual even-degree case of Theorem~\ref{descent} is not
excluded here. Nor does this note prove the existence of an
interior perfect split in a smaller wing, verify every required
leaf-kernel condition, or establish the stability bijection.
It supplies no whole-original, open-question or Lean credit.
\end{document}
